% Source: MQ20b_v1.html

\documentclass[11pt]{article}
\usepackage[margin=1in]{geometry}
\usepackage{amsmath}
\usepackage{amssymb}
\usepackage{siunitx}
\usepackage{enumitem}

\title{MQ20b: Kinetic Theory / Ideal Gas}
\author{}
\date{}

\begin{document}
\maketitle

\section*{MQ20b: Kinetic Theory / Ideal Gas}

\begin{enumerate}[leftmargin=*,label=\textbf{Question \arabic*.}]

\item \hfill \textit{1 pts}

A Carnot engine has an efficiency of 0.51. If the temperature of the hot reservoir is 542°C, calculate the temperature of the cold reservoir in Kelvin.

\item \hfill \textit{1 pts}

A Carnot engine ejects 1,067 J to the cold reservoir. If that reservoir has a temperature of 263 K and the engine does 467 J of work, calculate the temperature of the hot reservoir in Kelvin.

\item \hfill \textit{1 pts}

Suppose 19.1 kg of water at 0°C freezes to ice at 0°C. Calculate the water's entropy change in kJ/K. (\emph{L}$_{f,water}$= 3.35 × 10$^{5}$ J/kg)



Δ\emph{S} = \rule{2.5cm}{0.4pt} kJ/K

\item \hfill \textit{1 pts}

Gallium melts to a liquid at 29.76°C. Suppose 21 kg of liquid gallium at this temperature is placed in large room at 17°C. Calculate the change in entropy of the universe during the time that the gallium freezes solid, with neither gallium nor the room significantly changing temperature. \emph{L}$_{f,gallium}$ = 8.04 x 10$^{4}$ J/kg.

\item \hfill \textit{1 pts}

Calculate the change in entropy of water when 2.1 kg of water at 29°C warms to  54°C.   [\emph{c}$_{water}$ = 4186 J/(kg⋅K)]

\item \hfill \textit{1 pts}

A box consists of 9 equal-volume compartments, separated by partitions. Suppose 4 compartments contain a sum total of 2.5 moles of an ideal gas, with the other compartments completely evacuated. The partitions are removed so the gas freely expands, filling all compartments. Calculate the gas's entropy change during this process.

\item \hfill \textit{1 pts}

Suppose 18.8 kg of steam at 100°C cools to ice at 0°C. Calculate the change in entropy of the water in kJ/K.

[\emph{L}$_{v,water}$ = 2.26 × 10$^{6}$ J/kg, \emph{L}$_{f,water}$ = 3.35 × 10$^{5}$ J/kg, \emph{c}$_{water}$ = 4186 J/(kg⋅K)]



Δ\emph{S} = \rule{2.5cm}{0.4pt} kJ/K

\end{enumerate}

\end{document}
